\chapter{Fatigue}

\section*{Definition}
Fatigue is a subjective lack of physical or mental energy that is not relieved adequately by usual rest and may impair daily functioning. It differs from objective muscle weakness and from sleepiness.

\section*{When to See a Doctor}
See your doctor if:
\begin{itemize}
  \item Fatigue is persistent, worsening, or interferes with daily life
\end{itemize}

\section*{How It Develops}
Multifactorial symptom representing potential combination of peripheral muscle fatigue and central nervous system mechanisms across numerous conditions.

\section*{Causes}
\begin{itemize}
  \item Depression/anxiety disorders
  \item Anemia
  \item Obstructive sleep apnea
  \item Hypothyroidism
  \item Medication side effects
\end{itemize}

\section*{Related Symptoms}
\begin{itemize}
  \item Weakness/concentration difficulty
  \item Sleep disturbance
  \item Mood changes
  \item Malaise
  \item Reduced exercise tolerance
\end{itemize}

\section*{Medical Evaluation}
\begin{itemize}
  \item Your doctor will ask about your symptoms: when they started, what triggers them, and what makes them better or worse.
  \item They will review your medical history, medications, and any recent illnesses or exposures.
  \item A physical examination looks for signs of anemia, thyroid disease, infection, cardiopulmonary disease, mood disorder, sleep disorder, or other causes.
  \item Testing is guided by the history and examination and may include a blood count, metabolic or thyroid testing, or sleep assessment when indicated.
\end{itemize}

\section*{Treatment Options}
\begin{itemize}
  \item Treatment depends on the underlying cause identified during your evaluation.
  \item Treatment targets the identified cause; there is no general medication that safely treats unexplained fatigue.
  \item Lifestyle changes, rest, and home care measures can help you feel better.
  \item Your doctor will schedule follow-up visits to monitor your progress and adjust treatment as needed.
\end{itemize}


\section*{Self-Care and Home Management}
\begin{itemize}
  \item Prioritize adequate rest and sleep to support the body's healing and recovery processes.
  \item Maintain a balanced diet with regular meals and stay well-hydrated.
  \item Monitor symptoms and track any changes in severity, frequency, or associated features.
  \item Avoid overexertion and gradually increase activity levels as tolerated.
  \item Seek medical evaluation if symptoms persist, worsen, or interfere significantly with daily function.
\end{itemize}

\section*{References}
\begin{thebibliography}{99}
\bibitem{fatigue:rosenthal2008}
Rosenthal TC, Majeroni BA, Pretorius R. Fatigue: an overview. \textit{Am Fam Physician}. 2008;78(10):1167--1170.
\bibitem{fatigue:sharpe1998}
Sharpe M, Wilks D. Fatigue. \textit{BMJ}. 2002;325(7362):480--483.
\bibitem{fatigue:papadopoulos2012}
Papadopoulos FC, Sjoqvist PO. Fatigue. In: \textit{Harrisons Principles of Internal Medicine}. 18th ed. McGraw-Hill; 2012.
\bibitem{fatigue:afari2003}
Afari N, Buchwald D. Chronic fatigue syndrome: a review. \textit{Am J Psychiatry}. 2003;160(2):221--236.
\bibitem{fatigue:craig2003}
Craig T, Kakumanu S. Chronic fatigue syndrome: evaluation and treatment. \textit{Am Fam Physician}. 2002;65(6):1083--1090.
\bibitem{fatigue:prins2006}
Prins JB, van der Meer JW, Bleijenberg G. Chronic fatigue syndrome. \textit{Lancet}. 2006;367(9507):346--355.
\end{thebibliography}
